\chapter{Riemann 几何}

\begin{solution}{riemannian-conformal-metric}
$e^{2u(p)}>0$，所以 $\widetilde g_p=e^{2u(p)}g_p$ 仍为正定对称双线性型，
其系数随 $p$ 光滑。向量长度变为
$|v|_{\widetilde g}=e^u|v|_g$，故
\[
 L_{\widetilde g}(\gamma)=\int e^{u(\gamma(t))}|\dot\gamma(t)|_gdt.
\]
$n$ 维中 $\det(e^{2u}g_{ij})=e^{2nu}\det(g_{ij})$，所以
$dV_{\widetilde g}=e^{nu}dV_g$。
\end{solution}

\begin{solution}{riemannian-length-reparametrization}
设 $t=\phi(s)$，其中 $\phi$ 分段 $C^1$、单调递增并把新参数区间映到
$[a,b]$。链式法则给
\[
 \left|\frac d{ds}\gamma(\phi(s))\right|_g
 =|\dot\gamma(\phi(s))|_g\phi'(s).
\]
一元变量代换遂给新曲线长度等于
$\int_a^b|\dot\gamma(t)|_gdt$。若反向参数化，则出现 $|\phi'|$，长度仍
不变而定向改变。
\end{solution}

\begin{solution}{riemannian-energy-length}
在 $[a,b]$ 上对 $|\dot\gamma|$ 与常数函数 $1$ 使用
Cauchy--Schwarz：
\[
 L(\gamma)^2\le(b-a)\int_a^b|\dot\gamma|^2dt
 =2(b-a)E(\gamma).
\]
等号恰在 $|\dot\gamma(t)|$ 与 $1$ 线性相关，即速率几乎处处为常数时成立；
对光滑正则曲线就是处处常速。
\end{solution}

\begin{solution}{riemannian-volume-coordinate-change}
若 $x=x(y)$，Jacobian 矩阵为 $J=(\partial x^i/\partial y^a)$，则新坐标中
度量矩阵为 $g_y=J^{\mathsf T}g_xJ$。因此
\[
 \sqrt{\det g_y}=|\det J|\sqrt{\det g_x}.
\]
另一方面
\[
 |dx^1\cdots dx^n|
 =|\det J|\,|dy^1\cdots dy^n|.
\]
变量代换正好给出同一密度，故
$\sqrt{\det g_{ij}}\,|dx^1\cdots dx^n|$ 不依赖坐标。
定向情形去掉绝对值，并由定向相容坐标保证 Jacobian 为正。
\end{solution}

\begin{solution}{riemannian-euclidean-connection}
直角坐标中 $g_{ij}=\delta_{ij}$ 为常数，故所有 Christoffel 符号为零，
\[
 \nabla_XY=X(Y^k)\partial_k.
\]
测地线方程为 $\ddot x^k=0$，解是仿射直线 $x(t)=p+tv$。曲率坐标式中的
导数项和二次项全部为零，所以 $R=0$。
\end{solution}

\begin{solution}{riemannian-koszul-derivation}
度量相容性分别给
\begin{align*}
 Xg(Y,Z)&=g(\nabla_XY,Z)+g(Y,\nabla_XZ),\\
 Yg(Z,X)&=g(\nabla_YZ,X)+g(Z,\nabla_YX),\\
 Zg(X,Y)&=g(\nabla_ZX,Y)+g(X,\nabla_ZY).
\end{align*}
前两式相加并减第三式。再用
$\nabla_XY-\nabla_YX=[X,Y]$ 三次替换，使含
$\nabla_XY$ 的两项保留，其余配成 Lie 括号，得到
\[
\begin{aligned}
 2g(\nabla_XY,Z)={}&Xg(Y,Z)+Yg(Z,X)-Zg(X,Y)\\
 &-g(X,[Y,Z])+g(Y,[Z,X])+g(Z,[X,Y]).
\end{aligned}
\]
\end{solution}

\begin{solution}{riemannian-christoffel-formula}
在 Koszul 公式中取 $X=\partial_i,Y=\partial_j,Z=\partial_\ell$。坐标向量
场两两交换，所以括号项为零，得到
\[
 2g_{k\ell}\Gamma^k_{ij}
 =\partial_i g_{j\ell}+\partial_jg_{i\ell}-\partial_\ell g_{ij}.
\]
两边乘 $g^{m\ell}$，便有
\[
 \Gamma^m_{ij}=\frac12g^{m\ell}
 (\partial_i g_{j\ell}+\partial_jg_{i\ell}-\partial_\ell g_{ij}).
\]
对称性 $\Gamma^m_{ij}=\Gamma^m_{ji}$ 对应零挠率。
\end{solution}

\begin{solution}{riemannian-geodesic-speed}
沿测地线利用度量相容性，
\[
 \frac d{dt}|\dot\gamma|^2
 =2g\!\left(\frac{D\dot\gamma}{dt},\dot\gamma\right)=0.
\]
所以速度恒定。若 $t=as+b$ 且 $a\ne0$，重新参数化后仍满足
$D_s(d\gamma/ds)=0$，这类参数称为仿射参数。一般单调参数会多出
$(d^2t/ds^2)\dot\gamma$ 项，曲线的像仍可能相同，但新参数通常不满足
测地线方程。
\end{solution}

\begin{solution}{riemannian-divergence-coordinate}
先在一个定向坐标域中把密度暂时写成正的顶次形式
$dV_g=\rho\,dx^1\wedge\cdots\wedge dx^n$，其中
$\rho=\sqrt{|g|}$。Cartan 公式及顶次形式的闭性给
\[
 \mathcal L_XdV_g=d(\iota_XdV_g)
 =\partial_i(\rho X^i)\,dx^1\wedge\cdots\wedge dx^n.
\]
由定义 $\mathcal L_XdV_g=(\operatorname{div}X)dV_g$，比较系数便得
$\operatorname{div}X=\rho^{-1}\partial_i(\rho X^i)$。反向坐标变换同时改变
坐标顶次形式的符号和定向，而密度使用 Jacobian 的绝对值；所得标量公式不变，
故同一计算适用于非定向流形。
\end{solution}

\begin{solution}{riemannian-green-identity}
乘积法则给
\[
 \operatorname{div}(u\operatorname{grad}v)
 =g(\operatorname{grad}u,\operatorname{grad}v)+u\Delta_gv.
\]
向量场 $u\operatorname{grad}v$ 在至少一个函数紧支撑时具有紧支撑。密度形式
的散度定理使上式左侧的积分为零，移项即得
\[
 \int_Mu\Delta_gv\,dV_g
 =-\int_Mg(\operatorname{grad}u,\operatorname{grad}v)dV_g.
\]
\end{solution}

\begin{solution}{riemannian-first-variation}
记 $T=\partial_t\Gamma$、$V=\partial_s\Gamma$，以下用沿参数族的协变导数。
度量相容性给
\[
 \frac d{ds}E(\Gamma_s)\bigg|_{s=0}
 =\int_a^bg(\nabla_VT,T)dt.
\]
由 $D_sT=D_tV$，上式的 $\nabla_VT$ 理解为 $D_sT$，等于 $D_tV$。
再用乘积法则，
\[
 g(\nabla_TV,T)=\frac d{dt}g(V,T)-g(V,\nabla_TT).
\]
积分得到第一变分公式。固定端点变分使 $V(a)=V(b)=0$；若曲线对所有这类
变分为临界点，则
$\int g(V,\nabla_TT)=0$ 对所有紧支撑变分场成立。局部坐标中的变分基本引理
给 $\nabla_TT=0$。反向代入公式即可。
\end{solution}

\begin{solution}{riemannian-sphere-geodesics}
设 $\gamma(t)$ 是单位球面上的单位速大圆。作为 $\R^{n+1}$ 中曲线，
$\ddot\gamma=-\gamma$，完全沿球面的法向；球面 Levi--Civita 导数是环境
导数的切向分量，故 $D\dot\gamma/dt=0$。小圆的环境加速度除法向分量外通常
还有球面切向分量。具体地，小圆可写为
$\gamma(t)=a+r(\cos(t/r)e_1+\sin(t/r)e_2)$，其中
$a\perp e_1,e_2$、$r^2+|a|^2=1$、$0<|a|<1$。它为单位速曲线，且
\[
 \frac{D\dot\gamma}{dt}
 =\ddot\gamma+\gamma
 =\frac{a-|a|^2\gamma}{r^2}\ne0.
\]
若该向量为零，则 $a=|a|^2\gamma$，取长度会迫使 $|a|=1$，矛盾。
所以任何这样的圆弧都不是测地线；$a=0$ 的大圆才消去此切向分量。
\end{solution}

\begin{solution}{riemannian-curvature-tensorial}
例如把第一变量换为 $fX$。利用
$\nabla_{fX}=f\nabla_X$ 和
$[fX,Y]=f[X,Y]-Y(f)X$，
\begin{align*}
 R(fX,Y)Z
 &=f\nabla_X\nabla_YZ-\nabla_Y(f\nabla_XZ)
   -\nabla_{f[X,Y]-Y(f)X}Z\\
 &=fR(X,Y)Z,
\end{align*}
其中两个 $Y(f)\nabla_XZ$ 相消。第二变量由反对称性得到。第三变量换成
$fZ$ 时，二阶求导产生的 $X(f)$、$Y(f)$ 和二阶标量导数项与括号项逐一
相消，余下 $fR(X,Y)Z$。因此曲率在每点只依赖三个向量的点值。
\end{solution}

\begin{solution}{riemannian-jacobi-equation}
对测地线变分有 $D_tT=0$。沿 $s$ 求导，并用参数协变导数交换式和
$D_sT=D_tV$，得到
\[
 0=D_sD_tT=D_tD_sT+R(V,T)T=D_t^2V+R(V,T)T.
\]
这就是 Jacobi 方程。Euclidean 空间 $R=0$ 且协变导数是普通导数，所以
$J''(t)=0$，全部解为 $J(t)=A+tB$，其中 $A,B$ 为常向量。

若 $\gamma(t)=\exp_p(tv)$ 且 $J(0)=0$、$DJ/dt(0)=W$，考察测地线变分
$\Gamma(s,t)=\exp_p(t(v+sW))$。其变分场为
$d(\exp_p)_{tv}(tW)$，初值恰为 $(0,W)$；Jacobi 方程初值解唯一，故
$J(t)=d(\exp_p)_{tv}(tW)$。这同时说明 $d(\exp_p)_{t_0v}$ 有非零核等价于
沿 $\gamma$ 存在非零 Jacobi 场在 $0,t_0$ 同时为零。
\end{solution}

\begin{solution}{riemannian-incomplete-interval}
度量就是通常的 $dx^2$，两点间最短长度为 $|x-y|$，所以
$d_g(x,y)=|x-y|$。序列 $x_n=1/(n+2)$（$n\in\N$）是 Cauchy 列，却在 $(0,1)$ 中没有极限，
故不完备。测地线为 $x(t)=a+bt$；例如从 $1/2$ 以速度 $-1$ 出发的极大
测地线 $x(t)=1/2-t$ 的极大区间为 $(-1/2,1/2)$，有限时间到达缺失端点，
故测地线也不完备。
\end{solution}

\begin{solution}{riemannian-second-variation-sphere}
设 $n\geq2$。沿单位速大圆取平行正交规范标架
$T,E_2,\ldots,E_n$。单位球面曲率为 $1$，故
$R(J,T)T=J-g(J,T)T$。把 Jacobi 场写成
$J=aT+\sum_{j=2}^nb_jE_j$，方程化为
\[
 a''=0,\qquad b_j''+b_j=0.
\]
若 $J(0)=0$，则 $a(t)=ct$、$b_j(t)=c_j\sin t$。再要求
$J(\ell)=0$：当 $0<\ell<\pi$ 时，$c=c_j=0$，不存在非零两端为零的 Jacobi
场；当 $\ell=\pi$ 时，任取非零垂直平行场 $E$，都有
$J(t)=\sin t\,E(t)$。所以 $\pi$ 确为首次共轭时刻，而不只是一个发生退化的
时刻。对 $S^1$ 则只有切向分量 $aT$，两端为零迫使 $a=0$；没有非零
垂直方向，因而不存在共轭点。

二阶变分也给出同一现象。形如 $V=\phi E$、
$\phi(0)=\phi(\ell)=0$ 的变分场满足
\[
 I(V,V)=\int_0^\ell\bigl((\phi')^2-\phi^2\bigr)dt.
\]
所需的区间 Wirtinger 不等式可直接证明。令 $\lambda=\pi/\ell$、
$s(t)=\sin(\lambda t)$，把 $\phi=s\psi$ 代入并分部积分，得到
\[
 \int_0^\ell\left((\phi')^2-\lambda^2\phi^2\right)dt
 =\int_0^\ell s^2(\psi')^2dt\geq0.
\]
先在 $[\varepsilon,\ell-\varepsilon]$ 上计算；边界项为 $ss'\psi^2$。
由于 $\phi$ 光滑且在两端为零，$\psi=\phi/s$ 在端点有光滑延拓，所以边界项
趋于零，给出上述等式。若 $\ell<\pi$，此不等式给
$\int(\phi')^2\ge(\pi/\ell)^2\int\phi^2>\int\phi^2$，二阶变分正定。
在 $\ell=\pi$ 时取 $\phi(t)=\sin t$，等号成立，正好得到上述 Jacobi 场。
\end{solution}
